\chapter{F-Strings, Variable Annotations, and the Compact Dict (Python 3.6)}

Python 3.6 is the release that made modern Python \emph{feel} modern. Three changes did it:
\textbf{f-strings} (PEP 498) replaced \texttt{\%} and \texttt{.format()} with compile-time
interpolation; \textbf{variable annotations} (PEP 526) extended type hints from function signatures
to any binding; and the \textbf{compact dict} turned the core mapping type ordered and 30--40\%
smaller---an ``implementation detail'' in 3.6 that became a language guarantee in 3.7. We teach each
with \emph{current} mechanics, including the \textbf{PEP 701} f-string overhaul (3.12) that the
original 3.6 descriptions predate.

\section*{Section Index}
\begin{itemize}
  \item \textbf{12.1} Formatted string literals (PEP 498) and the PEP 701 overhaul
  \item \textbf{12.2} Variable annotations (PEP 526)
  \item \textbf{12.3} The compact, ordered dictionary
  \item \textbf{12.4} Performance and anti-patterns
  \item \textbf{12.5} Summary and cross-references
\end{itemize}

\section{Formatted string literals (PEP 498) and the PEP 701 overhaul}

\textbf{Why this exists.} The pre-3.6 options pay runtime cost: \texttt{\%}-formatting uses rigid C
printf routines and a temporary tuple; \texttt{"...".format(x)} does an attribute lookup, a
function call, \emph{and} re-parses the format string every call. \textbf{f-strings} move
interpolation to \textbf{compile time}: the parser splits the literal into constant chunks and
embedded expression AST sub-trees, and the compiler emits inline opcodes.

\begin{lstlisting}[language=Python, caption={An f-string compiles to inline formatting opcodes --- no .format() call.}]
import dis
def fstring_format(name, age):
    return f"Name: {name!r:>10}, Age: {age}"
dis.dis(fstring_format)
\end{lstlisting}

Verified output (CPython 3.13.5):

\begin{lstlisting}[caption={Output}]
  6   LOAD_CONST        1 ('Name: ')
      LOAD_FAST         0 (name)
      CONVERT_VALUE     2 (repr)        # the !r conversion
      LOAD_CONST        2 ('>10')
      FORMAT_WITH_SPEC                  # apply the :>10 format spec
      LOAD_CONST        3 (', Age: ')
      LOAD_FAST         1 (age)
      FORMAT_SIMPLE                     # no conversion, no spec
      BUILD_STRING      4               # C-level concat of the 4 pieces
      RETURN_VALUE
\end{lstlisting}

\textbf{This is modern (3.13) bytecode and differs from every pre-3.12 text.} The original 3.6
implementation used a single \texttt{FORMAT\_VALUE} opcode with packed flag bits. \textbf{PEP 701
(Python 3.12)} reimplemented f-strings in the PEG grammar and split formatting into
\texttt{CONVERT\_VALUE}, \texttt{FORMAT\_SIMPLE}, and \texttt{FORMAT\_WITH\_SPEC}, finished by
\texttt{BUILD\_STRING}. If you see \texttt{FORMAT\_VALUE}, you are reading pre-3.12 output.

\begin{lstlisting}[language=Python, caption={Format specs, conversions, the = debug form, and PEP 701 relaxations.}]
x = 3.14159
print("spec:", f"{x:.2f}")
print("conv:", f"{'hi'!r}")
print("debug:", f"{x=}")
val = {"k": "v"}
print("nested same-quotes (PEP 701, 3.12+):", f"{val["k"]}")
print("multiline expression (PEP 701):", f"{
    1 + 2
}")
\end{lstlisting}

Verified output (CPython 3.13.5):

\begin{lstlisting}[caption={Output}]
spec: 3.14
conv: 'hi'
debug: x=3.14159
nested same-quotes (PEP 701, 3.12+): v
multiline expression (PEP 701): 3
\end{lstlisting}

\textbf{PEP 701 lifted the old restrictions}: reusing the f-string's own quote character,
backslashes, multiline expressions, and comments inside the braces all now work. The trade-off
remains: f-strings format \emph{eagerly} (§12.4).

\section{Variable annotations (PEP 526)}

\textbf{Why this exists.} PEP 3107 allowed annotations on function parameters/returns; \textbf{PEP
526} extended the syntax to any binding---module, class, and local variables---giving static
checkers a place to read types and the typing ecosystem its foundation.

\textbf{The static boundary---annotations are metadata, not checks.} The interpreter never
type-checks an annotation. At module/class level the compiler emits \texttt{SETUP\_ANNOTATIONS} and
records each in \texttt{\_\_annotations\_\_}; at function-local level annotations are
\textbf{discarded entirely}---zero runtime cost.

\begin{lstlisting}[language=Python, caption={Class annotations are stored; function-local annotations are stripped.}]
import dis

class Profile:
    name: str = "Anon"
    age: int                      # annotation only, no value

print("Profile.__annotations__:", Profile.__annotations__)

def process():
    y: int = 42                   # the ': int' is discarded by the compiler
    return y
dis.dis(process)
\end{lstlisting}

Verified output (CPython 3.13.5):

\begin{lstlisting}[caption={Output}]
Profile.__annotations__: {'name': <class 'str'>, 'age': <class 'int'>}
 26   LOAD_CONST  1 (42)
      STORE_FAST  0 (y)
 27   LOAD_FAST   0 (y)
      RETURN_VALUE
\end{lstlisting}

\begin{lstlisting}[language=Python, caption={An unresolved forward reference raises at class-definition time; quote it to defer.}]
try:
    class Node:
        parent: Node              # Node is not bound yet during its own body
except NameError as e:
    print("NameError:", e)

class NodeOK:
    parent: "NodeOK"              # string forward reference -- stored unevaluated
print("string forward ref:", NodeOK.__annotations__)
\end{lstlisting}

Verified output (CPython 3.13.5):

\begin{lstlisting}[caption={Output}]
NameError: name 'Node' is not defined
string forward ref: {'parent': 'NodeOK'}
\end{lstlisting}

Frameworks resolve string annotations with \texttt{typing.get\_type\_hints()}. \textbf{PEP 563}
(\texttt{from \_\_future\_\_ import annotations}) and \textbf{PEP 649} (lazy evaluation, default in
3.14) are covered in Vol V; the type \emph{system} is Vol XV.

\section{The compact, ordered dictionary}

\textbf{Why this exists.} \texttt{dict} backs every namespace, object \texttt{\_\_dict\_\_}, and
kwargs bundle, so its layout matters. Before 3.6 a dict was \emph{one} sparse table of 24-byte
slots kept $\le 2/3$ full---many empty 24-byte slots wasting memory and scattering cache lines.

\textbf{The compact design (Raymond Hettinger).} 3.6 split the table: \textbf{\texttt{dk\_indices}}
(a small sparse array of integer indices, 1/2/4/8 bytes each) and \textbf{\texttt{dk\_entries}}
(a dense contiguous array of \texttt{(hash, key, value)} appended in insertion order). A lookup
hashes into \texttt{dk\_indices}, reads the small integer, and indexes \texttt{dk\_entries}. Empty
slots cost 1--8 bytes, cutting dict memory \textbf{$\sim$30--40\%}, and iteration order = insertion
order, for free.

\begin{lstlisting}[language=Python, caption={The compact dict preserves insertion order (impl detail in 3.6, guaranteed since 3.7).}]
d = {}
for k in ["z", "a", "m", "b"]:
    d[k] = 1
print("insertion order:", list(d))
\end{lstlisting}

Verified output (CPython 3.13.5):

\begin{lstlisting}[caption={Output}]
insertion order: ['z', 'a', 'm', 'b']
\end{lstlisting}

Order was an implementation detail in 3.6 and a \textbf{language guarantee in 3.7}. The
senior-engineer contrast: Java provides this only via a separate \texttt{LinkedHashMap}; in Python
the default \texttt{dict} is ordered, at no extra cost, because it falls out of the compact layout.

\section{Performance and anti-patterns}

\begin{itemize}
  \item \textbf{f-strings are fastest but format eagerly}---in logging use
    \texttt{logging.info("x=\%s", x)} (formatted only if enabled), not \texttt{logging.info(f"x=\{x\}")}.
  \item \textbf{Never f-string untrusted data into SQL/HTML/shell}---interpolation is not escaping.
  \item \textbf{Annotations are not runtime validation}---\texttt{x: int = "nope"} runs; use
    pydantic/validators (Vol XV).
  \item \textbf{Function-local annotations cost nothing; class/module ones cost a little}
    (\texttt{SETUP\_ANNOTATIONS} at import).
  \item \textbf{Relying on dict order is safe (3.7+)}---but sets are unordered.
\end{itemize}

\section{Summary and cross-references}

\begin{itemize}
  \item \textbf{f-strings} (PEP 498) interpolate at compile time (3.13:
    \texttt{CONVERT\_VALUE}/\texttt{FORMAT\_SIMPLE}/\texttt{FORMAT\_WITH\_SPEC}/\texttt{BUILD\_STRING});
    \textbf{PEP 701} removed nesting/multiline/backslash restrictions. Eager---avoid in logging.
  \item \textbf{Variable annotations} (PEP 526) are metadata, not checks: stored at module/class
    scope, stripped at function scope; forward references need quoting.
  \item The \textbf{compact dict} (3.6) cuts memory $\sim$30--40\% and makes \textbf{insertion
    order} intrinsic---a guarantee since 3.7.
\end{itemize}

\textbf{Cross-references.} \texttt{\%}/\texttt{.format} and the \texttt{string} module $\rightarrow$
Vol XI. Deferred annotations (PEP 563/649) $\rightarrow$ Vol V. The type system $\rightarrow$
Vol XV. Dataclasses $\rightarrow$ Chapter 13. CPython dict internals $\rightarrow$ Vol XII. f-string
codegen and the PEG grammar $\rightarrow$ Vol X.
